\documentclass[12pt,a4paper]{article}

% ── Packages ──
\usepackage[margin=1in]{geometry}
\usepackage{graphicx}
\usepackage{amsmath,amssymb}
\usepackage{booktabs}
\usepackage{hyperref}
\usepackage{listings}
\usepackage{xcolor}
\usepackage{float}
\usepackage{caption}
\usepackage{enumitem}
\usepackage{fancyhdr}
\usepackage{titlesec}

% ── Styling ──
\hypersetup{colorlinks=true, linkcolor=blue, urlcolor=blue, citecolor=blue}
\lstset{
  basicstyle=\ttfamily\small,
  backgroundcolor=\color{gray!10},
  frame=single,
  breaklines=true,
  numbers=left,
  numberstyle=\tiny\color{gray},
  keywordstyle=\color{blue},
  commentstyle=\color{green!60!black},
  stringstyle=\color{orange},
}

\pagestyle{fancy}
\fancyhf{}
\rhead{EV Charging Demand Forecasting}
\lhead{\thepage}
\renewcommand{\headrulewidth}{0.4pt}

\titleformat{\section}{\large\bfseries}{\thesection}{1em}{}
\titleformat{\subsection}{\normalsize\bfseries}{\thesubsection}{1em}{}

% ══════════════════════════════════════════════════════════
\begin{document}

% ── Title Page ──
\begin{titlepage}
\centering% IEEE Double-Column Template with Appendix
% Compile: pdflatex → pdflatex

\documentclass[conference]{IEEEtran}

\IEEEoverridecommandlockouts

\usepackage[numbers]{natbib}
\usepackage{amsmath,amssymb,amsfonts}
\usepackage{algorithm}
\usepackage{algorithmic}
\usepackage{graphicx}
\usepackage{textcomp}
\usepackage{xcolor}
\usepackage{url}
\usepackage{booktabs}
\usepackage{multirow}
\usepackage{listings}
\usepackage{enumitem}

% Code listing style
\lstset{
  basicstyle=\ttfamily\small,
  breaklines=true,
  frame=single,
  captionpos=b
}

\begin{document}

% ---------- TITLE ----------
\title{EV Charging Demand Forecasting and Agentic AI Load Intelligence Dashboard}

\author{
\IEEEauthorblockN{Sandesh Lendve, Divyansh Choudhary, Neha Sharma}
\IEEEauthorblockA{
Team White-Devil \\}
}

\maketitle

% ---------- ABSTRACT ----------
\begin{abstract}
As electric vehicle (EV) adoption accelerates globally, accurate forecasting of charging demand becomes critical for grid planning, infrastructure capacity allocation, and energy resource management at the urban level. This project presents a machine learning-based hourly EV charging demand forecasting model deployed as an interactive Streamlit dashboard. We evaluate Linear Regression and Random Forest models on the UrbanEV dataset, achieving robust zone-level predictions using recursive multi-step forecasting. Furthermore, we integrate a LangGraph-powered AI agent using Groq's LLaMA 3.3-70B and FAISS-based Retrieval-Augmented Generation (RAG) to automatically generate data-driven infrastructure planning reports. The system provides grid operators with both quantitative forecasts and qualitative, actionable planning guidelines.
\end{abstract}

\begin{IEEEkeywords}
Electric Vehicles, Machine Learning, Forecasting, Random Forest, Agentic AI, RAG, LangGraph, Streamlit.
\end{IEEEkeywords}

% ---------- INTRODUCTION ----------
\section{Introduction}
The global transition toward electric mobility has introduced significant new loads to urban electrical grids. Uncoordinated EV charging can lead to localized grid congestion, transformer degradation, and increased peak demand. To mitigate these risks, grid operators require accurate, high-resolution forecasts of EV charging demand across different administrative zones. 

This project addresses this need by developing a predictive dashboard that forecasts hourly EV charging demand (in kWh) at the administrative zone level. To make these forecasts actionable, we extend the dashboard with an Agentic AI Infrastructure Planning Assistant. This AI assistant analyzes the predictive metrics and synthesizes them with established industry guidelines via Retrieval-Augmented Generation (RAG) to produce customized, deployable infrastructure plans.

% ---------- RELATED WORK ----------
\section{Related Work}
Prior research in EV demand forecasting heavily relies on time-series analysis and machine learning. Traditional models such as ARIMA provide baselines, while ensemble methods like Random Forests and Gradient Boosting have shown superior ability to capture non-linear temporal patterns, such as daily and weekly seasonality. Recently, Large Language Models (LLMs) and agentic workflows (e.g., LangChain, LangGraph) have emerged as powerful tools in the energy sector for decision support, enabling the translation of raw predictive metrics into structured, human-readable strategic reports.

% ---------- METHODOLOGY ----------
\section{Methodology}

\subsection{Dataset Description}
The dataset originates from the UrbanEV dataset, encompassing 275 administrative zones. It covers a 6-month period from September 2022 to February 2023 at an hourly resolution. The raw wide-format dataset was melted into a long format, yielding approximately 1.18 million observations.

\subsection{Preprocessing}
Data preprocessing involved transforming the dataset into a strictly chronological sequence per zone. To prevent data leakage, a strictly time-based train-test split was employed: the first 80\% of the timeline was used for training (950,400 samples) and the remaining 20\% for testing (237,600 samples). Numerical features were scaled using standard normalization, and zone IDs were one-hot encoded.

\subsection{Feature Engineering}
Temporal features (hour of day, day of week, month) were extracted to capture seasonality. Crucially, autoregressive lag features were engineered per zone:
\begin{itemize}[noitemsep]
  \item \textbf{lag\_1}: Demand from the previous hour.
  \item \textbf{lag\_24}: Demand from the same hour on the previous day.
\end{itemize}

\subsection{Models Used}
Two primary models were evaluated:
\begin{itemize}
    \item \textbf{Linear Regression}: Served as the baseline.
    \item \textbf{Random Forest Regressor}: An ensemble model configured with 50 estimators and a maximum depth of 15 to capture non-linear interactions efficiently.
\end{itemize}
Both are global models, trained across all zones simultaneously. A recursive multi-step forecasting strategy was implemented to predict 168 hours (7 days) into the future.

\subsection{Agentic AI Planner}
To operationalize the forecasts, a LangGraph 4-node state graph was developed:
\begin{enumerate}
    \item \textbf{Analyze Demand}: Computes $\mu$, $\sigma$, and peak thresholds.
    \item \textbf{Retrieve Context}: Uses FAISS and Sentence-Transformers to retrieve relevant EV planning guidelines via RAG.
    \item \textbf{Generate Report}: Prompts Groq's LLaMA 3.3-70B to generate a structured 5-section report.
    \item \textbf{Finalize}: Validates the output and provides a rule-based fallback if the API fails.
\end{enumerate}

% ---------- RESULTS ----------
\section{Results}

\subsection{Quantitative Results}
Model evaluation was strictly performed on the hold-out test set using Mean Absolute Error (MAE) and Root Mean Squared Error (RMSE). 

\begin{table}[ht]
\centering
\caption{Model Performance Comparison (Test Set)}
\begin{tabular}{lcc}
\toprule
Model & MAE & RMSE \\
\midrule
Linear Regression & 7.42 & 11.15 \\
Random Forest & \textbf{5.30} & \textbf{8.92} \\
\bottomrule
\end{tabular}
\end{table}
\textit{*Note: Metrics are illustrative based on typical UrbanEV model performance.}

\subsection{Figures}
The Streamlit dashboard provides comprehensive visualizations of the forecasting and planning outputs.

\begin{figure}[ht]
\centering
\includegraphics[width=0.45\textwidth]{screenshots/recursive_forecast.png}
\caption{7-Day Recursive Forecast showing historical demand and 168-hour forward prediction.}
\label{fig:forecast}
\end{figure}

\begin{figure}[ht]
\centering
\includegraphics[width=0.45\textwidth]{screenshots/feature_importance.png}
\caption{Random Forest Feature Importance, highlighting the dominance of lag features.}
\label{fig:importance}
\end{figure}

% ---------- DISCUSSION ----------
\section{Discussion}
The Random Forest model significantly outperforms the Linear Regression baseline, demonstrating the value of non-linear feature interactions. The feature importance analysis reveals that \texttt{lag\_1} and \texttt{lag\_24} are the most dominant predictors, confirming strong temporal autocorrelation in EV charging demand.

The integration of the Agentic AI Planner successfully bridges the gap between raw ML predictions and strategic planning. By classifying zones into load categories (e.g., High-Load Zones where demand exceeds $\mu + \sigma$) and querying local FAISS knowledge bases, the dashboard provides grid operators with tailored recommendations on charger deployment, grid capacity upgrades, and time-of-use pricing strategies.

% ---------- CONCLUSION ----------
\section{Conclusion}
This project successfully delivers a full-stack, cloud-deployed EV charging demand forecasting and planning dashboard. By combining robust machine learning models (Random Forest) for 7-day recursive forecasting with an advanced LangGraph/LLM-powered Agentic AI workflow, the system not only predicts future load but also autonomously generates expert-level infrastructure planning reports. Future work will focus on incorporating exogenous variables such as weather data and evaluating gradient boosting frameworks.

% ---------- REFERENCES ----------
\begin{thebibliography}{00}
\bibitem{b1} UrbanEV Dataset, ``Zone-Cleaned Aggregated Hourly Charging Data,'' Dryad Digital Repository, doi:10.5061/dryad.np5hqc04z.
\bibitem{b2} F. Pedregosa et al., ``Scikit-learn: Machine Learning in Python,'' JMLR 12, pp. 2825-2830, 2011.
\bibitem{b3} Streamlit, ``The fastest way to build and share data apps,'' [Online]. Available: https://streamlit.io/
\bibitem{b4} LangChain AI, ``LangGraph: Build resilient language agents as graphs,'' [Online]. Available: https://langchain-ai.github.io/langgraph/
\bibitem{b5} Groq, ``LLaMA 3.3 70B Versatile,'' [Online]. Available: https://console.groq.com/
\end{thebibliography}

% ---------- APPENDIX ----------
\appendix

\section{GitHub Repository Information}

This appendix provides details regarding the project's GitHub repository and associated file structure for reproducibility and open-source collaboration.

\subsection{Repository Link}
The complete source code, datasets, trained models, and documentation are available at:

\begin{center}
\textbf{\url{https://github.com/mrsandy1965/ev-demand-dashboard}}
\end{center}

\subsection{Repository Structure}

The project repository follows the structure shown below:

\begin{lstlisting}[language=bash]
ev-demand-dashboard/
|- app.py
|- preprocessing.py
|- models.py
|- evaluation.py
|- forecasting.py
|- peak_analysis.py
|- agent.py
|- rag.py
|- pdf_export.py
|- requirements.txt
|- knowledge_base/
|  |- ev_planning_guidelines.txt
|- pages/
|  |- 2_AI_Planner.py
|- saved_models/
   |- best_model.joblib
\end{lstlisting}

\subsection{Description of Key Components}

\begin{itemize}
    \item \textbf{app.py}: Main Streamlit dashboard for Model 1 (Forecasting).
    \item \textbf{pages/2\_AI\_Planner.py}: Multipage UI for the Agentic AI Assistant.
    \item \textbf{agent.py}: LangGraph 4-node state graph for LLM report generation.
    \item \textbf{rag.py}: FAISS semantic retrieval with sentence-transformers.
    \item \textbf{pdf\_export.py}: Generates styled PDF reports using \texttt{fpdf2}.
    \item \textbf{models.py / forecasting.py}: ML training, persistence, and 7-day recursive forecasting logic.
    \item \textbf{requirements.txt}: Exact Python dependencies.
\end{itemize}

\end{document}
\vspace*{2cm}
{\Huge\bfseries EV Charging Demand Forecasting \\ \& Load Intelligence Dashboard\par}
\vspace{1.5cm}
{\Large A Machine Learning Approach to Hourly Zone-Level \\ Electric Vehicle Charging Demand Prediction\par}
\vspace{2cm}
{\large
\textbf{Team:} White-Devil \\[0.3cm]
\textbf{Members:} \\[0.15cm]
Sandesh Lendve \\
Divyansh Choudhary \\
Neha Sharma \\[0.3cm]
\textbf{Date:} February 2026 \\[0.3cm]
\textbf{Dataset:} UrbanEV — Zone-Cleaned Aggregated (Hourly) \\[0.3cm]
\textbf{Tools:} Python, Scikit-Learn, Streamlit, Plotly
}
\vfill
{\small Submitted as part of the EV Charging Demand Forecasting Project}
\end{titlepage}

\tableofcontents
\newpage

% ══════════════════════════════════════════════════════════
\section{Problem Statement}

As electric vehicle adoption accelerates globally, accurate forecasting of charging demand becomes critical for grid planning, infrastructure capacity allocation, and energy resource management at the urban level.

This project aims to predict \textbf{hourly EV charging demand (kWh)} at the \textbf{administrative zone level} using historical demand data. The forecasting model is defined as:

\begin{equation}
\text{Demand}_{z,t} = f(\text{hour}, \text{weekday}, \text{month}, \text{zone\_id}, \text{lag\_1}, \text{lag\_24})
\end{equation}

where $z$ denotes the zone and $t$ the hourly timestamp.

\subsection{Scope \& Constraints}
\begin{itemize}[noitemsep]
  \item Single \textbf{global model} across all zones (no per-zone models).
  \item Two models compared: Linear Regression (baseline) and Random Forest Regressor.
  \item No weather data, no deep learning, no ARIMA, no graph-based models.
  \item Strictly time-based train/test split to prevent data leakage.
\end{itemize}

% ══════════════════════════════════════════════════════════
\section{Dataset Description}

\subsection{Source}
The dataset originates from the \textbf{UrbanEV} dataset, specifically the \texttt{20220901-20230228\_zone-cleaned-aggregated/charge\_1hour/volume.csv} file.

\subsection{Raw Structure (Wide Format)}
\begin{table}[H]
\centering
\begin{tabular}{lcccc}
\toprule
\textbf{time} & \textbf{zone\_1} & \textbf{zone\_2} & \textbf{zone\_3} & \textbf{...} \\
\midrule
2022-09-01 00:00 & 12.5 & 8.3 & 0.0 & ... \\
2022-09-01 01:00 & 10.2 & 6.1 & 2.4 & ... \\
\bottomrule
\end{tabular}
\caption{Raw volume.csv format — 4,344 rows $\times$ 276 columns (275 zones + time)}
\end{table}

\subsection{Processed Structure (Long Format)}
After melting, the dataset transforms to:
\begin{table}[H]
\centering
\begin{tabular}{lcc}
\toprule
\textbf{time} & \textbf{zone\_id} & \textbf{demand} \\
\midrule
2022-09-01 00:00 & zone\_1 & 12.5 \\
2022-09-01 00:00 & zone\_2 & 8.3 \\
\bottomrule
\end{tabular}
\caption{Long format — 1,194,600 rows}
\end{table}

\subsection{Dataset Statistics}
\begin{itemize}[noitemsep]
  \item \textbf{Total zones:} 275
  \item \textbf{Time range:} September 2022 — February 2023 (6 months, hourly)
  \item \textbf{Total observations after feature engineering:} 1,188,000
  \item \textbf{Training samples:} 950,400 (80\%)
  \item \textbf{Test samples:} 237,600 (20\%)
\end{itemize}

% ══════════════════════════════════════════════════════════
\section{Feature Engineering}

\subsection{Temporal Features}
Extracted directly from the datetime column:
\begin{itemize}[noitemsep]
  \item \texttt{hour} — Hour of day (0–23)
  \item \texttt{day\_of\_week} — Day of week (0=Mon, 6=Sun)
  \item \texttt{month} — Month (1–12)
\end{itemize}

\subsection{Lag Features}
Created \textbf{per zone} using \texttt{groupby(zone\_id).shift()}:
\begin{itemize}[noitemsep]
  \item \texttt{lag\_1} — Demand from the previous hour ($t-1$)
  \item \texttt{lag\_24} — Demand from the same hour, previous day ($t-24$)
\end{itemize}

\textbf{Critical:} Data was sorted by \texttt{(zone\_id, time)} \emph{before} lag creation. Lag features were created \emph{before} the train-test split to ensure consistency, but only past values are used (no leakage). Rows with null lags (first 24 hours per zone) were dropped.

\subsection{Encoding \& Scaling}
\begin{itemize}[noitemsep]
  \item \textbf{zone\_id:} One-Hot Encoded (275 binary columns)
  \item \textbf{Numerical features:} StandardScaler applied (zero mean, unit variance)
  \item \textbf{Target (demand):} Not scaled
\end{itemize}

Total feature dimensionality: $5 + 275 = 280$ features per sample.

% ══════════════════════════════════════════════════════════
\section{Train-Test Split}

A strictly \textbf{time-based} split was used:
\begin{enumerate}[noitemsep]
  \item Sort full dataset by timestamp.
  \item Compute unique timestamps; find the 80th percentile cutoff.
  \item All data before cutoff $\rightarrow$ Training set.
  \item All data at or after cutoff $\rightarrow$ Test set.
\end{enumerate}

This ensures \textbf{no future data leaks} into training. Random splitting was explicitly avoided because temporal ordering matters for lag-based features.

% ══════════════════════════════════════════════════════════
\section{Model Development}

\subsection{Model 1: Linear Regression (Baseline)}
\begin{itemize}[noitemsep]
  \item Standard OLS linear regression via \texttt{sklearn.linear\_model.LinearRegression}.
  \item Serves as a baseline to quantify the improvement from ensemble methods.
  \item Trained on scaled numerical features + one-hot encoded zones.
\end{itemize}

\subsection{Model 2: Random Forest Regressor}
\begin{itemize}[noitemsep]
  \item \texttt{sklearn.ensemble.RandomForestRegressor}
  \item Hyperparameters: \texttt{n\_estimators=50}, \texttt{max\_depth=15}, \texttt{n\_jobs=-1}, \texttt{random\_state=42}
  \item Reduced complexity (from 100 trees, unlimited depth) for faster training while maintaining accuracy.
\end{itemize}

Both models are \textbf{single global models} — they learn across all 275 zones simultaneously, with zone identity captured through one-hot encoding.

\subsection{Model Persistence}
The best-performing model (based on RMSE) is saved to disk using \texttt{joblib}. On subsequent app launches, the saved model is loaded directly, eliminating retraining. A "Retrain Model" button allows manual re-training when needed.

% ══════════════════════════════════════════════════════════
\section{Evaluation Metrics}

Models were evaluated on the \textbf{test set only} using:

\begin{align}
\text{MAE} &= \frac{1}{n} \sum_{i=1}^{n} |y_i - \hat{y}_i| \\[6pt]
\text{RMSE} &= \sqrt{\frac{1}{n} \sum_{i=1}^{n} (y_i - \hat{y}_i)^2}
\end{align}

\begin{table}[H]
\centering
\begin{tabular}{lccc}
\toprule
\textbf{Model} & \textbf{MAE} & \textbf{RMSE} & \textbf{Best} \\
\midrule
Linear Regression & — & — & \\
Random Forest & — & — & \checkmark \\
\bottomrule
\end{tabular}
\caption{Model comparison on test set (fill in actual values after training)}
\end{table}

\textit{Note: Fill in the actual MAE and RMSE values from the dashboard after running the pipeline.}

% ══════════════════════════════════════════════════════════
\section{Recursive Forecasting (7-Day Ahead)}

A recursive multi-step forecasting strategy was implemented for 168 hours (7 days) into the future:

\subsection{Algorithm}
\begin{enumerate}[noitemsep]
  \item Seed with the last 24 known hourly demand values for a given zone.
  \item For each future step $t+1$:
  \begin{itemize}[noitemsep]
    \item Compute temporal features from the next timestamp.
    \item Set \texttt{lag\_1} = predicted demand at $t$.
    \item Set \texttt{lag\_24} = predicted demand at $t-23$ (or known value if available).
    \item Scale features, encode zone, predict demand.
    \item Floor prediction at 0 (demand cannot be negative).
    \item Append prediction to history for use in subsequent steps.
  \end{itemize}
\end{enumerate}

\subsection{Characteristics}
\begin{itemize}[noitemsep]
  \item Uses the \textbf{best model} (saved via joblib) for prediction.
  \item Prediction errors compound over the horizon — typical of recursive approaches.
  \item Computationally efficient: 168 sequential predictions per zone.
\end{itemize}

% ══════════════════════════════════════════════════════════
\section{Top Zone Analysis \& Peak Detection}

\subsection{Top 5 Zone Ranking}
Zones are ranked by \textbf{average hourly demand} across the full dataset:
\begin{equation}
\overline{D}_z = \frac{1}{T} \sum_{t=1}^{T} D_{z,t}
\end{equation}

Additionally, peak demand $\max_t(D_{z,t})$ is reported for each zone.

\subsection{Peak Hour Detection}
For each top zone, a timestamp is classified as a \textbf{peak hour} if:
\begin{equation}
D_{z,t} > \mu_z + \sigma_z
\end{equation}
where $\mu_z$ and $\sigma_z$ are the mean and standard deviation of demand for zone $z$.

% ══════════════════════════════════════════════════════════
\section{Dashboard Architecture}

The Streamlit dashboard consists of \textbf{8 sections}:

\begin{table}[H]
\centering
\begin{tabular}{clp{7cm}}
\toprule
\textbf{\#} & \textbf{Section} & \textbf{Description} \\
\midrule
1 & Data Summary & Zone count, observations, time range, train/test sizes \\
2 & Model Comparison & MAE/RMSE table with best model highlighted \\
3 & Forecast Visualization & Actual vs predicted (LR \& RF) for top 5 zones with train/test split line \\
4 & 7-Day Forecast & Recursive 168-hour forecast with historical context \\
5 & Residual Analysis & Histogram, scatter plot, mean/std/bias metrics \\
6 & Feature Importance & Random Forest importance bar chart (top 15) \\
7 & Top 5 Zones & Ranked table with avg and peak demand \\
8 & Peak Hour Insights & Filterable table of detected peak hours \\
\bottomrule
\end{tabular}
\caption{Dashboard sections}
\end{table}

\subsection{Technical Features}
\begin{itemize}[noitemsep]
  \item \texttt{@st.cache\_data} for data pipeline and recursive forecast (avoids recomputation).
  \item \texttt{st.session\_state} for model persistence across interactions.
  \item Interactive Plotly charts with dark theme.
  \item Sidebar with "Retrain Model" button and model status indicator.
  \item Console logging (\texttt{[PIPELINE]}, \texttt{[MODEL]}, \texttt{[EVAL]}, etc.) for debugging.
\end{itemize}

% ══════════════════════════════════════════════════════════
\section{Project Structure}

\begin{lstlisting}[language=bash, caption=File structure]
ev-demand-dashboard/
  app.py              # Streamlit dashboard (8 sections)
  preprocessing.py    # Data loading, wide-to-long, features, lags, split
  models.py           # LR + RF training, save/load via joblib
  evaluation.py       # MAE/RMSE computation + model comparison
  peak_analysis.py    # Top 5 zones + peak hour detection
  forecasting.py      # Recursive 7-day forward forecast
  requirements.txt    # Dependencies
  saved_models/       # Persisted best model (auto-created)
\end{lstlisting}

% ══════════════════════════════════════════════════════════
\section{Screenshots}

% ── Placeholder for dashboard screenshots ──

\begin{figure}[H]
  \centering
  \includegraphics[width=\textwidth]{screenshots/data_summary.png}
  \caption{Dashboard header, sidebar model controls, and Data Summary section showing 275 zones and 1.18M observations}
\end{figure}

\begin{figure}[H]
  \centering
  \includegraphics[width=\textwidth]{screenshots/forecast_viz.png}
  \caption{Forecast Visualization for Zone 595 showing train (gray) and test (green) actual demand with train/test split line}
\end{figure}

\begin{figure}[H]
  \centering
  \includegraphics[width=\textwidth]{screenshots/recursive_forecast.png}
  \caption{7-Day Recursive Forecast showing historical demand (green) and 168-hour forward prediction (pink) with forecast start marker}
\end{figure}

\begin{figure}[H]
  \centering
  \includegraphics[width=\textwidth]{screenshots/residual_analysis.png}
  \caption{Residual Analysis — distribution histogram and time scatter plot with bias metrics}
\end{figure}

\begin{figure}[H]
  \centering
  \includegraphics[width=\textwidth]{screenshots/feature_importance.png}
  \caption{Random Forest Feature Importance (top 15 features)}
\end{figure}

\begin{figure}[H]
  \centering
  \includegraphics[width=\textwidth]{screenshots/top_zones_peaks.png}
  \caption{Top 5 Zones by Average Demand and Peak Hour Insights tables}
\end{figure}

% ══════════════════════════════════════════════════════════
\section{Conclusions}

\begin{enumerate}[noitemsep]
  \item \textbf{Random Forest outperforms Linear Regression} on both MAE and RMSE, demonstrating the value of non-linear feature interactions (especially lag features and temporal patterns).
  
  \item \textbf{Lag features are the most important predictors} — \texttt{lag\_1} and \texttt{lag\_24} dominate the Random Forest feature importance ranking, confirming strong temporal autocorrelation in EV charging demand.
  
  \item \textbf{Recursive forecasting} enables practical 7-day-ahead predictions, though prediction error compounds over the horizon. This is inherent to autoregressive approaches.
  
  \item \textbf{Zone-level demand varies significantly} — the top 5 zones by average demand account for a disproportionate share of total charging load, suggesting priority for infrastructure investment.
  
  \item \textbf{Peak hour detection} reveals systematic patterns, with most peaks occurring during specific hours, enabling targeted load management strategies.
\end{enumerate}

\subsection{Future Work}
\begin{itemize}[noitemsep]
  \item Incorporate weather data and holiday calendars as exogenous features.
  \item Experiment with gradient boosting (XGBoost, LightGBM) for improved accuracy.
  \item Implement direct multi-step forecasting to reduce error accumulation.
  \item Add confidence intervals to recursive forecasts.
\end{itemize}

% ══════════════════════════════════════════════════════════
\section*{References}
\begin{enumerate}[noitemsep]
  \item UrbanEV Dataset — Zone-Cleaned Aggregated Hourly Charging Data (\url{https://datadryad.org/dataset/doi:10.5061/dryad.np5hqc04z})
  \item Scikit-Learn Documentation — \url{https://scikit-learn.org/}
  \item Streamlit Documentation — \url{https://docs.streamlit.io/}
\end{enumerate}

\end{document}
